\documentclass[11pt]{article}

\usepackage[margin=30mm]{geometry}
\usepackage[T1]{fontenc}
\usepackage{mathptmx}
\usepackage{amsmath,amssymb,amsthm,mathtools}
\usepackage{booktabs,tabularx}
\usepackage[colorlinks=true,linkcolor=blue,urlcolor=blue]{hyperref}

\newcommand{\R}{\mathbb{R}}
\newcommand{\diff}{\mathrm{d}}
\newcommand{\cL}{\mathcal{L}}
\newcommand{\Rea}{\operatorname{Re}}

\theoremstyle{plain}
\newtheorem{proposition}{Proposition}
\newtheorem{lemma}{Lemma}

\theoremstyle{remark}
\newtheorem{remark}{Remark}

\title{A Step-by-Step Repair of the Killed-Resolvent Proof\\
\large Why the old boundary argument fails and why the corrected one works}
\author{}
\date{\today}

\begin{document}
\maketitle

\begin{abstract}
The current paper reaches the right-looking $3\times3$ linear system by using
one function in two incompatible ways.  It first treats the resolvent as zero
below the liquidation barrier, but then applies a full-line exponential
integral to its particular and homogeneous pieces.  This note explains the
problem using the constant particular solution $1/q$, introduces the required
notation from scratch, and gives a corrected proof.  The repair is simple in
principle: a downward jump integral must stop at the barrier.  For an
exponential term, this truncated integral equals the familiar rational term
minus one extra exponential term.  The last two rows of the $3\times3$ system
make these extra terms cancel; the first row imposes absorption at the
barrier.  The numerical matrix can therefore be retained, but the old
boundary interpretation must be replaced.
\end{abstract}

\section{What the proof is trying to calculate}

Fix a moment order $k$ and a time-Laplace variable $q$.  The object of interest
is the Laplace transform of a payoff that is set to zero after the log health
factor crosses its liquidation boundary.  In the paper's original coordinate,
the log-ratio process is denoted by $z$ and liquidation occurs at
\[
a:=-h_0.
\]
The resolvent is needed only for $z>a$, but the generator contains jump
integrals whose landing points can lie below $a$.  This is where the boundary
treatment matters.

It is clearer to measure distance from liquidation.  Define
\begin{equation}\label{eq:x-definition}
x:=z-a=z+h_0.
\end{equation}
Then $x>0$ means that the position is alive, $x=0$ is the barrier, and $x<0$
means that a downward jump has crossed the barrier.  Let
\[
u(q,x)
\]
denote the killed resolvent in this translated coordinate.  Absorption means
\begin{equation}\label{eq:killed-extension}
u(q,x)=0,\qquad x\leq0.
\end{equation}

\section{Why the current proof does not work}

Consider one downward exponential jump phase with rate $r>0$.  If jump size is
$y>0$, then a jump from $x$ lands at $x-y$.  Applied to the killed function,
the downward jump operator is
\begin{equation}\label{eq:old-killed-op}
K_{r,-}^{D}u(q,x)
:=\int_0^\infty u(q,x-y)r e^{-ry}\,\diff y
=\int_0^x u(q,x-y)r e^{-ry}\,\diff y.
\end{equation}
The second equality follows from \eqref{eq:killed-extension}: terms with
$y\geq x$ land at or below zero and contribute nothing.  Hence
\begin{equation}\label{eq:true-boundary-limit}
\lim_{x\downarrow0}K_{r,-}^{D}u(q,x)=0.
\end{equation}

The old proof then writes the interior formula as
\[
u(q,x)=P(q,x)+H(q,x),\qquad x>0,
\]
where $P$ is a particular solution and $H$ is a sum of characteristic
exponentials.  It uses linearity to replace
\eqref{eq:true-boundary-limit} by
\[
K_{r,-}[P](0)+K_{r,-}[H](0)=0,
\]
and evaluates these two terms using full-line exponential formulas.  That step
is not valid: $P$ and $H$ have only been introduced as pieces of an
\emph{interior} formula.  Their exponential formulas extend them to $x<0$,
whereas the killed function in \eqref{eq:killed-extension} is zero.

\subsection*{The contradiction is visible for the constant $1/q$}

For the survival probability, the particular solution inside the alive region
is the constant
\[
P(q,x)=\frac1q,\qquad x>0.
\]
There are now two different calculations:
\begin{align*}
\text{killed extension:}\quad
K_{r,-}^{D}[P](0^+)&=0,\\
\text{constant extended to the whole line:}\quad
\int_0^\infty \frac1q r e^{-ry}\,\diff y&=\frac1q.
\end{align*}
Both calculations are individually correct, but they concern different
functions below the barrier.  The old proof uses the first interpretation to
state the boundary condition and the second interpretation to compute its
components.  Linearity cannot connect two different extensions.

There is a second problem.  For a characteristic term $e^{\gamma x}$, the
full-line calculation requires
\[
\int_0^\infty e^{-\gamma y}r e^{-ry}\,\diff y<\infty,
\quad\text{or equivalently}\quad \Rea(\gamma)>-r.
\]
Some negative characteristic roots can lie to the left of the pole $-r$.
Then the integral diverges even though the rational expression
$r/(r+\gamma)$ is finite.  A proof based on the literal full-line integral
therefore cannot cover all roots used by the matrix.

The failure is thus not that the final matrix necessarily has the wrong
numbers.  The failure is that the stated boundary condition does not imply
those matrix rows by the reasoning given.

\section{Notation for the corrected proof}

The main symbols are collected here before the calculation begins.
\begin{center}
\centering
\renewcommand{\arraystretch}{1.25}
\begin{tabularx}{\linewidth}{lX}
\toprule
Symbol & Meaning \\
\midrule
$z$ & Original log-ratio state used in the paper. \\
$a=-h_0$ & Liquidation boundary in the $z$ coordinate. \\
$x=z+h_0$ & Distance from the liquidation boundary; the alive region is $x>0$. \\
$u(q,x)$ & Killed resolvent, set equal to zero for $x\leq0$. \\
$v(q,x)$ & Exponential formula proposed only on the interior $x>0$. \\
$r_{d,-}^{(k)}$ & Rate of downward exponential phase $d\in\{1,2\}$ under the $k$-tilted measure. \\
$\lambda_{d,-}^{(k)}$ & Intensity of downward phase $d$; it is positive. \\
$K_{r,-}^{D}$ & True killed downward operator; its integral stops at jump size $x$. \\
$\cL^{(k),D}$ & True generator of the process killed at $x=0$. \\
$\cL^{(k),\mathrm{rat}}$ & Algebraic ``rational'' operator used to identify particular solutions and characteristic roots. \\
$C_r[v]$ & Coefficient of the extra $e^{-rx}$ term created by truncating a downward integral. \\
\bottomrule
\end{tabularx}
\end{center}

\section{The true killed generator}

Fix a tilt order $k$.  Under the $k$-tilted measure, the log-ratio process has
two upward exponential jump phases and two downward exponential jump phases.
Index upward phases by $a\in\{1,2\}$ and downward phases by
$d\in\{1,2\}$.  Write their positive intensities as
$\lambda_{a,+}^{(k)}$ and $\lambda_{d,-}^{(k)}$, and their positive rates as
$r_{a,+}^{(k)}$ and $r_{d,-}^{(k)}$.  The exact intensity formulas are not
needed for the boundary argument; positivity and the two distinct downward
rates are what matter.

More explicitly, phase 1 comes from jumps of asset 1 and phase 2 from jumps of
asset 2.  Since $Z=X_1-X_2$, a negative jump of asset 2 is an upward jump of
$Z$, while a positive jump of asset 2 is a downward jump of $Z$.  The rates
and intensities are
\begin{align*}
r_{1,+}^{(k)}&=\frac1{\eta_{1,+}},
&\lambda_{1,+}^{(k)}&=\lambda_1p_1,\\
r_{2,+}^{(k)}&=\frac1{\eta_{2,-}}+k,
&\lambda_{2,+}^{(k)}&=\frac{\lambda_2(1-p_2)}{1+k\eta_{2,-}},\\
r_{1,-}^{(k)}&=\frac1{\eta_{1,-}},
&\lambda_{1,-}^{(k)}&=\lambda_1(1-p_1),\\
r_{2,-}^{(k)}&=\frac1{\eta_{2,+}}-k,
&\lambda_{2,-}^{(k)}&=\frac{\lambda_2p_2}{1-k\eta_{2,+}}.
\end{align*}
The condition $k\eta_{2,+}<1$ makes the last rate and intensity positive.

For a function defined on $x>0$, the generator of the killed process is
\begin{align}
\cL^{(k),D}f(x)
={}&\mu^{(k)}f'(x)+\frac{\sigma_Z^2}{2}f''(x)
 +\sum_{a=1}^2\lambda_{a,+}^{(k)}
       \bigl(K_{a,+}^{(k)}f(x)-f(x)\bigr) \notag\\
&+\sum_{d=1}^2\lambda_{d,-}^{(k)}
       \bigl(K_{d,-}^{(k),D}f(x)-f(x)\bigr),
\qquad x>0.                                      \label{eq:killed-generator}
\end{align}
Here all phase intensities are positive, and
\begin{align}
K_{a,+}^{(k)}f(x)
&:=\int_0^\infty f(x+y)r_{a,+}^{(k)}e^{-r_{a,+}^{(k)}y}\,\diff y,
                                                        \label{eq:up-op}\\
K_{d,-}^{(k),D}f(x)
&:=\int_0^x f(x-y)r_{d,-}^{(k)}e^{-r_{d,-}^{(k)}y}\,\diff y.
                                                        \label{eq:down-killed}
\end{align}
The upper limit $x$ in \eqref{eq:down-killed} is essential: a downward jump
larger than $x$ crosses the liquidation barrier, where the payoff is zero.

The drift $\mu^{(k)}$ and diffusion variance $\sigma_Z^2$ are the corresponding
tilted quantities.  In particular, $\sigma_Z>0$ means that absorption requires
the ordinary boundary condition $u(q,0)=0$.

The paper's rational characteristic function $\psi_Z^{(k)}(\beta)$ is useful,
but it must be interpreted carefully.  It is obtained by replacing each phase
operator acting on $e^{\beta x}$ by a rational multiplier.  It is an algebraic
device for finding exponentials that solve the interior equation.  It is not
the result of applying the killed operator at the boundary.

\section{The missing truncation identity}

\begin{lemma}[Truncated downward phase]\label{lem:truncated}
Let $r>0$ and let $\beta\in\mathbb{C}\setminus\{-r\}$.  On $x>0$,
\begin{equation}\label{eq:truncated-exp}
\int_0^x e^{\beta(x-y)}r e^{-ry}\,\diff y
=\frac{r}{r+\beta}e^{\beta x}
 -\frac{r}{r+\beta}e^{-rx}.
\end{equation}
The identity holds for every such $\beta$; no condition
$\Rea(\beta)>-r$ is required because the integral is over the finite interval
$[0,x]$.
\end{lemma}

\begin{proof}
Direct integration gives
\[
e^{\beta x}r\int_0^x e^{-(r+\beta)y}\,\diff y
=e^{\beta x}\frac{r}{r+\beta}
  \bigl(1-e^{-(r+\beta)x}\bigr),
\]
which is \eqref{eq:truncated-exp}.
\end{proof}

For a finite exponential sum
\[
v(x)=\sum_{\beta\in\mathcal E}a_\beta e^{\beta x},
\]
define its \emph{truncation-correction coefficient}
\begin{equation}\label{eq:phase-functional}
C_r[v]
:=\sum_{\beta\in\mathcal E}a_\beta\frac{r}{r+\beta}.
\end{equation}
This is simply a number calculated from the coefficients of the interior
formula $v$.  It is not the value of a jump integral at the boundary.
Lemma~\ref{lem:truncated} gives
\begin{equation}\label{eq:phase-decomposition}
K_{r,-}^{D}v(x)
=\sum_{\beta\in\mathcal E}
  a_\beta\frac{r}{r+\beta}e^{\beta x}
 -C_r[v]e^{-rx}.
\end{equation}
The first term is the expression produced by the rational calculation used in
the paper.  The second term is what that calculation omits when jumps crossing
the barrier are killed.  This omitted term is the key to the repair.

For example, if $v(x)=1/q$, then $C_r[v]=1/q$ and
\begin{equation}\label{eq:constant-truncated}
K_{r,-}^{D}v(x)=\frac1q-\frac1q e^{-rx}.
\end{equation}
This tends to zero as $x\downarrow0$, as it must.  The old proof kept the first
$1/q$ but did not explicitly derive the cancelling second term.

\section{The proposed repair}

The proposal has four steps.
\begin{enumerate}
\item Work only on the alive half-line $x>0$.
\item Use the rational equation to construct an interior candidate $v(q,x)$.
\item Apply the \emph{true truncated} downward operators to this candidate.
This produces two extra terms, one proportional to
$e^{-r_{1,-}^{(k)}x}$ and one proportional to
$e^{-r_{2,-}^{(k)}x}$.
\item Choose the three homogeneous coefficients so that the value at the
barrier is zero and both extra terms vanish.
\end{enumerate}

The next proposition carries out these steps.

Let $k\geq1$ be the integer moment order, let $b\in(0,1)$ be the liquidation
threshold parameter, and let $q$ be the time-Laplace variable.  Translation
removes $h_0$ from the forcing:
\[
G_k(x):=(e^x-b)^k
=\sum_{j=0}^k d_{k,j}e^{jx},
\qquad
d_{k,j}:=\binom{k}{j}(-b)^{k-j}.
\]
Assume that all exponentials in the forcing are in the positive-moment
domain.  In the bivariate Kou model this requires, in particular,
\begin{equation}\label{eq:moment-domain}
k\eta_{1,+}<1,
\qquad
k\eta_{2,+}<1.
\end{equation}
The second condition makes the $k$-tilt well defined; the first ensures that
the long-leg exponential moments and $\psi_Z^{(k)}(j)$ for $j\leq k$ are
finite.

Let $q$ initially lie to the right of the time-Laplace abscissa of convergence
and suppose $q\neq\psi_Z^{(k)}(j)$ for $j=0,\ldots,k$.  Define the algebraic
particular solution
\begin{equation}\label{eq:particular}
P_k(q,x):=
\sum_{j=0}^k
\frac{d_{k,j}}{q-\psi_Z^{(k)}(j)}e^{jx}.
\end{equation}
Suppose the rational characteristic equation
\[
\psi_Z^{(k)}(\gamma)=q
\]
has three simple roots $\gamma_4,\gamma_5,\gamma_6$ with negative real part,
none equal to a phase pole.  Consider
\begin{equation}\label{eq:interior-ansatz}
v_k(q,x):=P_k(q,x)+\sum_{m=4}^6B_m^{(k)}(q)e^{\gamma_m(q)x},
\qquad x>0.
\end{equation}

\begin{proposition}[Correct barrier conditions]\label{prop:correct}
Assume $\sigma_Z>0$ and
$r_{1,-}^{(k)}\neq r_{2,-}^{(k)}$.  The function
\begin{equation}\label{eq:killed-candidate}
\widehat U_k(q,x):=
\begin{cases}
v_k(q,x),&x>0,\\
0,&x\leq0,
\end{cases}
\end{equation}
solves
\begin{equation}\label{eq:killed-resolvent}
(q-\cL^{(k),D})\widehat U_k(q,x)=G_k(x),
\qquad x>0,
\end{equation}
with an absorbing boundary if and only if
\begin{equation}\label{eq:three-conditions}
v_k(q,0)=0,
\qquad
C_{r_{1,-}^{(k)}}[v_k]=0,
\qquad
C_{r_{2,-}^{(k)}}[v_k]=0.
\end{equation}
These conditions give the $3\times3$ system
\begin{equation}\label{eq:matrix-system}
\begin{pmatrix}
1&1&1\\[5pt]
\dfrac{r_{1,-}^{(k)}}{r_{1,-}^{(k)}+\gamma_4}&
\dfrac{r_{1,-}^{(k)}}{r_{1,-}^{(k)}+\gamma_5}&
\dfrac{r_{1,-}^{(k)}}{r_{1,-}^{(k)}+\gamma_6}\\[11pt]
\dfrac{r_{2,-}^{(k)}}{r_{2,-}^{(k)}+\gamma_4}&
\dfrac{r_{2,-}^{(k)}}{r_{2,-}^{(k)}+\gamma_5}&
\dfrac{r_{2,-}^{(k)}}{r_{2,-}^{(k)}+\gamma_6}
\end{pmatrix}
\begin{pmatrix}B_4^{(k)}\\B_5^{(k)}\\B_6^{(k)}\end{pmatrix}
=-
\begin{pmatrix}
P_k(q,0)\\[3pt]
C_{r_{1,-}^{(k)}}[P_k]\\[3pt]
C_{r_{2,-}^{(k)}}[P_k]
\end{pmatrix},
\end{equation}
where
\begin{equation}\label{eq:particular-functional}
C_r[P_k]
=\sum_{j=0}^k
\frac{d_{k,j}}{q-\psi_Z^{(k)}(j)}\frac{r}{r+j}.
\end{equation}
Thus \eqref{eq:matrix-system} is algebraically the same system as in the
paper, after substituting $x=z+h_0$.
\end{proposition}

\begin{proof}
Apply the killed generator \eqref{eq:killed-generator} to the interior ansatz
\eqref{eq:interior-ansatz}.  Upward phase integrals are not truncated by the
lower barrier.  For the positive forcing exponents they converge under
\eqref{eq:moment-domain}; for the negative characteristic roots they converge
automatically.  Their action is therefore the usual rational-symbol action.

For each downward phase, use \eqref{eq:phase-decomposition}.  Define
$\cL^{(k),\mathrm{rat}}$ to be the algebraic operator that keeps the first
term in \eqref{eq:phase-decomposition} but not its truncation correction.  It
is exactly the operator represented by the rational function
$\psi_Z^{(k)}$.  For $x>0$,
\begin{equation}\label{eq:operator-correction}
\cL^{(k),D}v_k(x)
=\cL^{(k),\mathrm{rat}}v_k(x)
-\sum_{d=1}^2\lambda_{d,-}^{(k)}
  C_{r_{d,-}^{(k)}}[v_k]e^{-r_{d,-}^{(k)}x}.
\end{equation}
By construction of $P_k$ and the characteristic roots,
\[
(q-\cL^{(k),\mathrm{rat}})v_k(x)=G_k(x).
\]
It follows from \eqref{eq:operator-correction} that
\begin{align}
(q-\cL^{(k),D})v_k(x)
=G_k(x)
 +\sum_{d=1}^2\lambda_{d,-}^{(k)}
  C_{r_{d,-}^{(k)}}[v_k]e^{-r_{d,-}^{(k)}x}.
                                                        \label{eq:residual}
\end{align}
Because the two rates are distinct and their intensities are positive, the
two extra exponentials are linearly independent.  They cannot cancel each
other for every $x>0$.  Equation \eqref{eq:killed-resolvent} therefore holds
exactly if and only if each correction coefficient is zero:
\[
C_{r_{1,-}^{(k)}}[v_k]=C_{r_{2,-}^{(k)}}[v_k]=0.
\]

Since $\sigma_Z>0$, the lower boundary is regular for the Brownian component
and killing imposes the Dirichlet condition $v_k(q,0)=0$.  These are the three
conditions in \eqref{eq:three-conditions}.  Applying them to
\eqref{eq:interior-ansatz} gives \eqref{eq:matrix-system} and
\eqref{eq:particular-functional}.  Standard uniqueness of the killed
resolvent in the relevant exponential-growth class then identifies the
candidate with the probabilistic resolvent.
\end{proof}

\section{Survival probability}

For survival, set $k=0$ and $G_0\equiv1$.  The algebraic particular solution
is $P_0=1/q$, hence
\[
P_0(0)=\frac1q,
\qquad
C_r[P_0]=\frac1q.
\]
The right-hand side of the survival system is therefore
$-(1/q,1/q,1/q)^\top$, exactly as in the paper.  The correct explanation is
not that the downward convolution of the killed resolvent at the barrier is
$1/q$; it is that the constant \emph{interior particular solution} produces a
barrier-truncation residual of size $(1/q)e^{-rx}$, which the homogeneous
terms must cancel.

\section{Why the proposal works while the old argument does not}

There are five concrete differences.
\begin{enumerate}
\item \emph{Below the barrier.}  The current proof uses the zero extension
when stating the boundary condition, but full-line exponential extensions when
evaluating its pieces.  The corrected proof never extends the interior
candidate below the barrier.

\item \emph{The constant particular solution.}  The current proof combines a
zero killed-boundary integral with the full-line identity $K[1/q]=1/q$.  The
corrected proof instead derives
\[
K^D[1/q](x)=q^{-1}(1-e^{-rx}),
\]
which tends to zero at the barrier without contradiction.

\item \emph{Roots beyond a phase pole.}  The full-line integral can diverge
when $\Rea(\gamma)\leq-r$.  The corrected integral is over the finite interval
$[0,x]$ and is valid for every $\gamma\neq-r$.

\item \emph{Meaning of the matrix.}  The current proof presents the last two
rows as jump-operator values at the boundary.  The corrected proof shows that
they cancel the coefficients of the two omitted functions
$e^{-r_{1,-}^{(k)}x}$ and $e^{-r_{2,-}^{(k)}x}$.

\item \emph{Direct verification.}  The corrected proof substitutes the
candidate into the true generator with truncated downward jumps.  The result
is the desired forcing plus two explicit error terms, and the matrix makes
both errors zero.
\end{enumerate}

The corrected proof works for a concrete reason that can be checked line by
line.  Start with the interior candidate $v_k$.  The rational construction
guarantees
\[
(q-\cL^{(k),\mathrm{rat}})v_k=G_k.
\]
The true killed generator differs from the rational calculation only by the
two truncation terms displayed in \eqref{eq:operator-correction}.  Therefore
substitution into the true equation gives
\[
(q-\cL^{(k),D})v_k
=G_k
+\lambda_{1,-}^{(k)}C_{r_{1,-}^{(k)}}[v_k]e^{-r_{1,-}^{(k)}x}
+\lambda_{2,-}^{(k)}C_{r_{2,-}^{(k)}}[v_k]e^{-r_{2,-}^{(k)}x}.
\]
The second and third matrix rows set the two $C$ coefficients to zero, so the
candidate satisfies the true killed PIDE on every $x>0$.  The first row gives
$v_k(q,0)=0$, so it also satisfies absorption.  Finally, uniqueness of the
killed resolvent identifies this verified candidate with the probabilistic
expectation.  Unlike the old proof, no step changes the function being used
below the barrier.

\section{Translation back to the paper's notation}

The original paper works in $z$ and writes
\[
c_{k,j}(h_0)e^{jz}
=\binom{k}{j}(-b)^{k-j}e^{jh_0}e^{jz}.
\]
At $x=z+h_0$, this becomes simply $d_{k,j}e^{jx}$.  Evaluation at the initial
state $z=0$ corresponds to $x=h_0$, so
\begin{equation}\label{eq:evaluate-initial}
\widehat U_k(q,0;h_0)
=P_k(q,h_0)+\sum_{m=4}^6B_m^{(k)}(q)e^{\gamma_m(q)h_0}.
\end{equation}
This also shows that, after translation, the barrier coefficients $B_m^{(k)}$
do not depend on $h_0$; the initial buffer enters only through evaluation at
$x=h_0$.

\section{Replacement guidance for the manuscript}

The revised paper should make the following changes.
\begin{enumerate}
\item Translate the barrier to zero with $x=z+h_0$ before deriving the
resolvent.
\item Define the killed downward operator with upper limit $x$, as in
\eqref{eq:down-killed}.
\item Replace the claim
``$K_{d,-}[\widehat U](0^+)=0$ and linearity gives conditions on the
particular and homogeneous pieces'' by Lemma~\ref{lem:truncated} and the
residual calculation \eqref{eq:residual}.
\item Call $C_r[v]$ a truncation-correction coefficient.  Do not identify it
with the value of the killed jump operator at the boundary.
\item State the additional long-leg moment condition $k\eta_{1,+}<1$.
\item Initially derive the time-Laplace transform for $q$ to the right of its
abscissa of convergence.  Treat continuation to the complex Talbot contour as
a separate numerical-analytic step.
\end{enumerate}

\begin{remark}[Coincident downward rates]
If $r_{1,-}^{(k)}=r_{2,-}^{(k)}$, the two downward exponential kernels combine
into a single phase.  The two residual functions in \eqref{eq:residual} are
then identical, so one should reduce the rational symbol and solve the
lower-dimensional system.  Describing this case only as a singular
$3\times3$ matrix obscures the probabilistic reason for the degeneracy.
\end{remark}

\section{Conclusion}

The original $3\times3$ coefficients can survive the correction, but the
boundary proof cannot.  The matrix rows involving
$r/(r+\gamma)$ are cancellation conditions for barrier-truncation residuals,
not values of the killed phase operator at the barrier.  This distinction
removes the extension inconsistency and remains valid for characteristic
roots lying beyond a downward phase pole, where the full-line convolution
interpretation is unavailable.

\end{document}
